\section{Framework} \label{sec:32-research01-framework}
\subsection{Problem statement}
\begin{foldable}
  Given a black-box pre-trained teacher network $T$ and a small set of labeled images $\mathcal{D}=\{x_{i},y_{i}\}_{i=1}^{N}$, our goal is to train a student network $S$ on $\mathcal{D}$ such that $S$ can approximate the performance of $T$.
\end{foldable}

\begin{foldable}
  A direct solution is to use standard \ac{KD}~\cite{hinton2015distilling}, which trains $S$ to match both the ground-truth labels $y_{i}$ and the teacher's predictions $T(x_{i})$ via minimizing:
  \begin{equation}
    \mathcal{L}_\text{KD} = \mathbb{E}_{(x_{i}, y_{i}) \sim \mathcal{D}} \Bigl[{
        (1 - \lambda) \mathcal{L}_{\text{CE}}(S(x_{i}), y_{i}) +
        \lambda \mathcal{L}_\text{KL}(S(x_{i}), T(x_{i}))
    }\Bigr]
    \label{eq:StandardKD}
  \end{equation}
  where $\mathcal{L}_{\text{CE}}$ and $\mathcal{L}_\text{KL}$ are the cross-entropy and the Kullback-Leibler divergence losses, $S(x_{i})$ and $T(x_{i})$ are the student's and teacher's predictions (\ie, class probabilities), and $\lambda$ is a coefficient to balance the losses.
  Note that the \textit{temperature} factor in the original paper is not used as it requires access to the teacher's logits, which violates our black-box teacher assumption.
  Moreover, with a small distillation set $\mathcal{D}$, common \ac{KD} approaches lead to a \textit{sub-optimal} model as they require lots of training images~\cite{hinton2015distilling,kim2018paraphrasing,ahn2019variational,tian2020contrastive}.
\end{foldable}

\begin{foldable}
  \textbf{Proposed Method.}
  To perform black-box few-shot \ac{KD}, we propose our novel method \textbf{DivBFKD}.
  Our method has two phases: Generation and Distillation.
  In the Generation phase, we train a WGAN to produce synthetic images.
  In the Distillation phase, we train the student with the combined real and synthetic images.
\end{foldable}

\subsection{Generation phase.}
\begin{foldable}
  We train a WGAN model~\cite{gulrajani2017improved} with several key changes to generate diverse synthetic images.
  Like standard WGAN, our WGAN consists of a generator $G$ and a discriminator $D$.
  Iteratively, $G$ produces a synthetic image $\tilde{x} = G(z)$ from a random latent vector $z \sim \mathcal{N}(0, \mathbf{I})$.
  Meanwhile, $D$ predicts a scalar score indicating the likelihood of an input image being real.
  Jointly trained with opposite objectives, $G$ tries to fool $D$ by producing realistic images while $D$ tries to best distinguish between real images $x$ and synthetic images $\tilde{x}$.
  These objectives can be formulated into the generator loss $\mathcal{L}_G$ and the discriminator loss $\mathcal{L}_{D}$:
  \begin{equation}
    \mathcal{L}_G = -\mathbb{E}_{z \sim \mathcal{N}(0, \mathbf{I})} [D(G(z))],\label{eq:loss-G-wgan}
  \end{equation}
  \begin{equation}
    \mathcal{L}_{D} = \mathbb{E}_{z \sim \mathcal{N}(0, \mathbf{I})} [D(G(z))] - \mathbb{E}_{x \sim \mathcal{D}}[D(x)].\label{eq:loss-D-wgan}
  \end{equation}
\end{foldable}

\begin{foldable}
  However, if we only use the few-shot set $\mathcal{D}$ to train the WGAN, $G$ will generate synthetic images similar to the images in $\mathcal{D}$~\cite{karras2020training,yang2021data}, which are not diverse enough to facilitate effective distillation.
  To improve diversity of images generated by $G$, we propose the use of synthetic images that the teacher $T$ can predict confidently.
  We refer to these images as \textit{high-confidence} images.
\end{foldable}

\begin{foldable}
  \textbf{High-confidence images.}
  Given a synthetically generated (fake) image $\tilde{x}$, we define its \textit{confidence-score} $c_{\tilde{x}} = \max{k \in \{0, \ldots, K-1\}}T(\tilde{x})$, by taking the maximum over $T$'s predictive probability vector.
  For example, given three classes $\{0, 1, 2\}$ and class probabilities $T(\tilde{x})=[0.1, 0.7, 0.2]$, then $c_{\tilde{x}}=0.7$.
  We define $\tilde{x}$ a high-confidence image if $c_{\tilde{x}}$ is above a confidence-threshold $\tau\in[0,1]$.
  Intuitively, when $c_{\tilde{x}}$ is high, $\tilde{x}$ is likely close to the teacher's training images as it can be confidently classified.
\end{foldable}

\begin{foldable}
  \textbf{Adaptive thresholds.} The critical factor in determining a high-confidence image is the confidence-threshold $\tau$.
  We observe that the teacher has a class-specific bias, \ie, it predicts different classes with different confidence levels.
  Take the dataset FMNIST~\cite{xiao2017fashion} for example, the averaged confidence-score is $>0.99$ for classes `trouser' and `bag', but is only 0.75 for class `shirt' (Fig.~\ref{fig:03-framework-pre_clamp_tau}).
  In this case, fixing a single value $\tau=0.95$ may lead to high-confidence images being too lenient for `trouser' and `bags', but too strict for the class `shirt'.
  Thus, a single $\tau$ for all classes will make high-confidence images imbalanced.
  Further, considering any single class, the confidence-scores have a long-tailed distribution, where they concentrate near 1 and decay with lower values.
  Even though $\tau$ can be set to the averaged confidence-score, it would provide little flexibility and interpretability and also make high-confidence images biased towards \textit{outliers}.
\end{foldable}

\begin{figure}[ht]
  \centering
  \includegraphics[viewport=0bp 20bp 455bp 174bp,clip,width=0.75\columnwidth]{./figures/30-research01/33-research01-framework-preclamptau.pdf}
  \caption[DivBFKD confidence-score distributions on FMNIST.]{
    DivBFKD confidence-score distributions across classes on FMNIST.
    They concentrate around 1 and have a long-tail.
    Red bars show adaptive confidence-thresholds $\tau$.
  }
  \label{fig:03-framework-pre_clamp_tau}
\end{figure}

\begin{foldable}
  To address these problems, we use adaptive thresholds $\{\tau^{k}\}_{k=0}^{K-1}$ that is class-specific based on the $q$-quantile of the confidence-scores of real images.
  Specifically, for each of the $K$ classes:
  \begin{equation}
    \tau^{k} = \text{quantile}_{q}(\{c_{x_{i}} \mid(x_{i}, y_{i}) \in \mathcal{D} \text{ and } y_{i} = k\}_{i=1}^{N}), q\in[0,1].
    \label{eq:adaptive-threshold}
  \end{equation}
  We choose $q = 0.1$, \ie, high-confidence images must have teacher's confidence higher than at least 10\% of those for the real ones, as perceived by the teacher $T$.
  This way, $\tau$ adaptively adjusts to any dataset and any class-wise bias from the teacher.
  Fig.~\ref{fig:03-framework-pre_clamp_tau} illustrates $\tau$ as red bars using $q = 0.1$ on FMNIST.
\end{foldable}

\begin{foldable}
  \textbf{WGAN architecture.} We employ a deep convolutional GAN where the generator and discriminator have symmetrical architectures.
  The generator accepts latent vectors of dimension 100 for MNIST, 200 for FMNIST, and 256 for other datasets.
  The latent inputs are projected to 256 base feature maps of dimension 8$\times$8 (except for Imagenette, 7$\times$7), followed by batch normalization.
  These feature maps pass through upscaling convolutional blocks, each consisting of a nearest neighbor (2$\times$) upsample, a convolutional layer with half the channels, batch normalization, and leaky ReLU activation.
  Once the feature maps match the training image dimensions, they are compressed to grayscale or RGB via a final convolutional layer with sigmoid activation, mapping pixel values to $[0, 1]$.

  The discriminator processes input images through convolutional blocks, each with (0.5$\times$) downscaling convolution, batch normalization (except the first block), and leaky ReLU activation.
  The final block contains 256 base feature maps of dimension 8$\times$8 (except Imagenette, 7$\times$7), with each previous block having half the channels.
  These feature maps are passed through a linear layer to produce a scalar score.
\end{foldable}

\begin{foldable}
  \textbf{Our WGAN training scheme.}
  We aim to use the teacher's guidance to promote diversity of synthetic images.
  Fundamentally, synthetic images would be more diverse if the WGAN can learn on images that resemble the teacher's \textit{unknown} training images $\mathcal{D}^{*}$, whose diversity is superior to few-shot images $\mathcal{D}$.
  While there is no direct access to $\mathcal{D}^{*}$, high-confidence images serve as a satisfactory alternative for both quantity (new high-confidence images are generated at every WGAN training step) and quality (they likely reside close to $\mathcal{D}^{*}$, as we previously discussed).
  In other words, high-confidence images implicitly allow our WGAN to learn from $\mathcal{D}^{*}$.
\end{foldable}

\begin{figure}
  \centering
  \includegraphics[width=0.99\columnwidth]{./figures/30-research01/33-research01-framework-teaser.pdf}
  \caption[Illustration of our method DivBFKD.]{
    Illustration of our method DivBFKD.
    \textbf{(a)} \textbf{Generation:} We train our WGAN with the losses $\mathcal{L}_G$ and $\mathcal{L}_{D}^{\text{new}}$.
    When optimizing $G$, we construct the high-confidence set $\mathcal{H}_{n}=\{\tilde{x} = G(z) \mid c_{\tilde{x}} \geq \tau^{\tilde{y}}\}$ from synthetic images $\tilde{x}$, with their confidence-score $c_{\tilde{x}}$, pseudo-label $\tilde{y}$, and adaptive threshold $\tau^{\tilde{y}}$.
    When optimizing $D$, we sample $x$ from the combined real and high-confidence images set $\mathcal{D} \cup \mathcal{H}_{n-1}$.
    \textbf{(b) Distillation:} We augment few images in $\mathcal{D}$ with synthetic images in $\mathcal{D}_G$ to train the student $S$ on $\mathcal{D} \cup \mathcal{D}_G$.
  }
  \label{fig:03-framework-teaser}
\end{figure}

\begin{foldable}
  By this motivation, we introduce high-confidence images in the role of real images to the adversarial learning (Fig.~\ref{fig:03-framework-teaser}(a)), exposing our WGAN to a wider variety of training images.
  At each $n$-th training step:
  \begin{itemize}
    \item {
        The generator $G$ is trained with Eq~\ref{eq:loss-G-wgan}.
        For a fake image $\tilde{x}$, we compute its corresponding confidence-score $c_{\tilde{x}}$ and pseudo-label $\tilde{y} = \argmax_{k \in \{0, \ldots, K-1\}} T(\tilde{x})$.
        Using the adaptive thresholds $\{\tau^{k}\}_{k=0}^{K-1}$ computed with Eq~\ref{eq:adaptive-threshold}, we add qualified $\tilde{x}$ to the \textit{high-confidence} set $\mathcal{H}_{n} = \{\tilde{x} \mid c_{\tilde{x}} \geq \tau^{\tilde{y}}\}$.
        Then, $\mathcal{H}_{n}$ is stored to train the discriminator $D$ in the next step.
      }
    \item {
        Training the discriminator $D$ requires both real and fake images.
        However, our `real' images include real images that actually come from $\mathcal{D}$ and high-confidence images that come from the high-confidence set $\mathcal{H}_{n-1}$.
        Note that $\mathcal{H}_{n-1}$ is constructed in the previous step, and $\mathcal{H}_{0}=\emptyset$.
        With these changes, we train our discriminator with the new loss as:
      }
  \end{itemize}
  \begin{equation}
    \mathcal{L}_{D}^{\text{new}} = {
      \mathbb{E}_{z \sim \mathcal{N}(0, \mathbf{I})}[D(G(z))] -
      \mathbb{E}_{x \sim (\mathcal{D} \cup \mathcal{H}_{n-1})}[D(x)]
    }.
    \label{eq:loss-D-divbfkd}
  \end{equation}
\end{foldable}

\begin{foldable}
  Viewing as an adversarial game, when $D$ learns to appreciate high-confidence synthetic images as realistic, it also encourages $G$ to create more synthetic images like those.
  This interaction forms a positive feedback loop that helps $G$ to be more diverse.
  Furthermore, we emphasize that forming $\mathcal{H}_{n}$ is a \textit{cheap} operator and requires small memory.
  These modifications not only effectively improve the diversity of synthetic images, but also incur only trivial computation overhead.
\end{foldable}

\begin{foldable}
  \textbf{Summary of our novelty in WGAN.}
  We introduce a new training scheme for WGAN using high-confidence images based on the teacher-supervised adaptive thresholds.
  Since we aim to generate images close to the teacher's training images, our synthetic images are more \textit{diverse} and better cover \textit{unseen images}, significantly improving the student's generalization, as shown in \S\ref{sec:33-research01-experiments}.
\end{foldable}

\begin{algorithm}[ht]
  \SetKwInOut{KwIn}{Input}
  \SetKwInOut{KwOut}{Output}
  \SetKwComment{Comment}{}{}

  \caption{Diverse Black-box Few-Shot Distillation (\textbf{DivBFKD}).}
  \label{alg:algorithm-divbfkd}

  \KwIn{teacher $T$; labeled set $\mathcal{D} = \{x_{i}, y_{i}\}_{i=1}^{N}$; quantile $q$; synthetic budget $M$; batch size $B$}
  \KwOut{student network $S$}

  \vspace{-0.5\baselineskip}
  \Comment{\noindent\hrulefill}

  \vspace{0.5\baselineskip}
  Initialize generator $G$, discriminator $D$, student $S$, and high-confidence set $\mathcal{H}_0 \leftarrow \emptyset$ \\
  \Comment{$\triangleright$ \textbf{\textit{Stage 1: Generation}}}
  Infer $T$ on $\mathcal{D}$ to compute adaptive thresholds $\{\tau^k\}_{k=1}^K$ with $q$ \hfill Eq.~\ref{eq:adaptive-threshold} \\
  \While{$G$ not converged}{
    Sample latent vectors $\{z_i\}_{i=1}^B \sim \mathcal{N}(0,I)$ \\
    Generate synthetic images $\{\tilde{x}_i\}_{i=1}^B$ with $G$ \\
    Sample $\{x_i\}_{i=1}^B \sim (\mathcal{D} \cup \mathcal{H}_n)$ as real images \\
    Update $G$ on $\left\{\tilde{x}_i\right\}_{i=1}^B$ via $\mathcal{L}_G$ \hfill Eq.~\ref{eq:loss-G-wgan} \\
    Update $D$ on $\left\{x_i\right\}_{i=1}^B$ and $\left\{\tilde{x}_i\right\}_{i=1}^B$ via $\mathcal{L}_{D}^{new}$ \hfill Eq.~\ref{eq:loss-D-divbfkd} \\
    Compute pseudo-labels $\{\tilde{y}_i\}$ and confidence-scores $\{c_{\tilde{x}_i}\}$ for $\{\tilde{x}_i\}$ \\
    Construct $\mathcal{H}_n=\left\{\tilde{x_{i}} \mid c_{\tilde{x}_i}\ge\tau^{\tilde{y}_i}\right\}_{i=1}^B$ \\
  }

  \vspace{0.5\baselineskip}
  \Comment{$\triangleright$ \textbf{\textit{Stage 2: Distillation}}}
  Construct $\mathcal{D}_\text{KD}$ with $N$ real and $M$ synthetic images \hfill Eq.~\ref{eq:distill-set} \\
  \While{$S$ not converged}{
    Update $S$ on $\mathcal{D}_\text{KD}$ via $\mathcal{L}_\text{KD}$ \hfill Eq.~\ref{eq:loss-S}
  }

  \vspace{0.5\baselineskip}
  \Return{$S$}
\end{algorithm}

\subsection{Distillation phase}
\begin{foldable}
  After training WGAN in the Generation phase, we use it to generate synthetic images.
  We employ our generator $G$ to build a set of synthetic images $\mathcal{D}_G$, and their pseudo-labels (\ie, class probabilities) are obtained via the teacher $T$.
  Following BBKD~\cite{wang2020neural} and FS-BBT~\cite{nguyen2022black}, we construct a distillation set $\mathcal{D}_\text{KD}$ comprising of $N$ real images from $\mathcal{D}$ and $M$ synthetic images from $\mathcal{D}_G$:
  \begin{equation}
    \mathcal{D}_\text{KD} = {
      \{x_{i} \in \mathcal{D}\}_{i=1}^{N} \cup
      \{G(z_{j}) | z_{j} \sim \mathcal{N}(0, \mathbf{I})\}_{j=1}^{M}
    },
    \label{eq:distill-set}
  \end{equation}
  where $x_{i} \in \mathcal{D}$ and $\tilde{x}_{j} \in \mathcal{D}_G$ are real and synthetic images.
  Like other few-shot \ac{KD} methods~\cite{wang2020neural,nguyen2022black}, we train the student $S$ with $\mathcal{D}_\text{KD}$ using a cross-entropy loss:
  \begin{equation}
    \mathcal{L}_\text{KD} = \mathbb{E}_{x_{i} \sim \mathcal{D}_\text{KD}} \mathcal{L}_\text{CE}(S(x_{i}), T(x_{i})).
    \label{eq:loss-S}
  \end{equation}
\end{foldable}

\begin{foldable}
  \textbf{Class balancing.} Inspired by~\cite{chen2019data,zhang2022ideal}, we generate a balanced set of synthetic images.
  Recalling that $M$ is the budget of synthetic images for \ac{KD} and $K$ is the number of classes, each class would have $M / K$ synthetic images.
  We ensure class balance through \textit{rejection sampling} until we achieve the sufficient amount.
\end{foldable}

We summarize our method in Algorithm \ref{alg:algorithm-divbfkd}.
